% =====================================================================
% DISCUSSION SECTION — draft 1
% Paper #1, privacy_hdt_wesad
% Register: formal academic, British English. Interpretive, not a results recap.
% Grounds every claim in a specific result; distinguishes what is established
% from what is conjectured.
% =====================================================================

\section{Discussion}
\label{sec:discussion}

The results of Section~\ref{sec:results} are usually reported as a single
privacy--utility number at a single operating point. Read together, they show
that this figure conflates three separable phenomena. We take them in turn, then
draw out the consequence for practice and for the human digital-twin setting that
motivates the wider programme.

\subsection{Federation and privacy are distinct harms with distinct incidence}
\label{sec:disc-decomposition}

The penalty of federating the data and the penalty of privatising the model are
not merely additive contributions to a single utility loss; they fall on
different classes, for different reasons, and are separable only because the
matched-configuration design holds every other component fixed between stages.

Federation, introduced in isolation at Stage~\ref{sec:results-federation},
leaves aggregate utility statistically unchanged yet imposes its entire cost on
one class: amusement declines by $0.095$ while the other three classes are
unaffected. The mechanism is not a generic loss of capacity but the averaging of
model updates across participants, which attenuates precisely the
subject-specific signal on which amusement recognition depends. Amusement is the
class whose physiological expression is most strongly idiosyncratic; it is,
in the precise sense, the class that transfers least well between people, and it
is exactly that class the averaging harms.

User-level privacy, introduced at Stage~\ref{sec:results-user-dp}, behaves
differently. Its degradation is approximately uniform across classes; once the
floor and ceiling effects that confound naive retention ratios are accounted
for, no class-specific privacy penalty is detectable. The a~priori expectation
that differential privacy would fall hardest on the minority class is, in this
four-class setting, not borne out: amusement is already near its floor after
federation, so it has little further to fall, and the remaining classes decline
together.

Sample-level privacy, introduced at Stage~\ref{sec:results-sample-dp}, exhibits
the class-specific effect that user-level privacy does not. It eliminates the
rarest class entirely --- amusement $F_1$ is exactly zero across the mid-range of
the budget --- while retaining more aggregate utility than user-level privacy at
every budget. The cause is localised by the noise-free diagnostic to the additive
noise rather than the gradient clipping, and it follows directly from the unit of
protection: because sample-level privacy protects the individual window, and the
rarest class contributes only about four windows to a mini-batch of thirty-two,
the class's gradient signal is the first to be overwhelmed as noise is added.

The three phenomena therefore have three distinct incidences: federation harms
the class that is hard to transfer between people; user-level privacy harms all
classes approximately equally; sample-level privacy harms the class that is rare
in absolute count. The widely repeated statement that differential privacy harms
minority classes is, against these results, imprecise in two ways. It does not
separate the penalty of distributing the data from the penalty of the privacy
mechanism, and it does not distinguish the two privacy granularities, which
differ in whether they exhibit a class-specific penalty at all.

\subsection{The choice of granularity is a fairness decision, not a privacy one}
\label{sec:disc-fairness}

The comparison of the two granularities is the paper's central practical claim,
and it turns on a coincidence that the empirical-privacy measurement of
Section~\ref{sec:results-mia} makes visible. On aggregate utility the two
mechanisms diverge sharply: sample-level privacy dominates user-level privacy at
every budget, by as much as $0.30$ macro-$F_1$. On empirical privacy they are
indistinguishable: for both mechanisms, and at every budget from
$\varepsilon = 0.5$ to $16$, a calibrated membership-inference attack is reduced
to chance. The two independent non-private runs agree to within $0.0007$ AUC,
which bounds the measurement noise well below the effect being asserted absent.

If the two mechanisms deliver the same empirical protection, then the decision
between them cannot be justified on privacy grounds. What separates them is
where their utility cost falls. Sample-level privacy buys its superior aggregate
performance by abandoning the rarest class; user-level privacy spreads a larger
cost across all classes while preserving their relative standing. The choice is
thus a choice about which distribution of error is acceptable --- a fairness
decision --- wearing the appearance of a privacy--utility trade-off.

The practical hazard is concrete. A practitioner selecting a mechanism from
published aggregate scores would prefer sample-level privacy at every budget, and
would thereby adopt the configuration that silently removes the minority class.
In an affective-computing context the minority class is amusement; in a clinical
monitoring context the analogous minority is frequently the event of interest ---
the rare arrhythmia, the infrequent fall, the uncommon seizure --- and its
elimination would not be visible in any aggregate metric. We therefore argue that
per-class performance is not optional supplementary reporting in private
federated learning but a necessary condition for a defensible mechanism choice.

\subsection{Formal budgets are conservative against an attacker this model does
not admit}
\label{sec:disc-formal-empirical}

Across both mechanisms the formally weak guarantee at $\varepsilon = 8$ delivers
the same empirical protection as the formally strong guarantee at
$\varepsilon = 0.5$. This decoupling of the formal budget from the measured
attack success should be read with care, and specifically not as evidence that
differential privacy is unnecessary.

The correct reading is that the empirical leakage of the undefended model is
already low. The non-private attack AUC is only $0.549$, and the
difficulty-calibrated attack does not exceed the raw attack, indicating that the
model does not encode strong per-window membership signal for a more powerful
attack to exploit. The within-subject memorisation gap tells the same story: it
is small but reliable on the undefended model and contracts monotonically towards
zero as the budget tightens. The formal guarantee bounds a worst-case adversary
against a worst-case model; the model trained here, being small and
well-regularised under early stopping, is far from that worst case, and so the
formal budget overstates the risk it actually carries. This is consistent with
prior findings, outside the wearable-health setting, that well-generalising
models leak little membership information irrespective of formal guarantees
\cite{[verify: generalisation-vs-DP, arXiv 2110.05524]}, and with the theoretical
result that the membership advantage is bounded by $e^{\varepsilon}-1$
\cite{[verify: Yeom 2018]}. What the present study adds is that the decoupling
holds across both privacy granularities on the same pipeline, so it is a property
of this model class in this domain rather than of a particular mechanism.

The claim is bounded in an important way. A low attack AUC is evidence that this
attack, the strongest we could mount without shadow models that the cohort size
forbids, does not succeed; it is not proof that no attack could. We therefore
frame the finding as a decoupling of formal budget from \emph{measured} risk, not
as a licence to weaken formal guarantees, and we return to this in the
limitations.

\subsection{Implications for privacy-preserving human digital twins}
\label{sec:disc-hdt}

The wider programme this study serves concerns human digital twins: live,
model-based representations of an individual's physiological state. A digital
twin whose unit of representation is the person requires a privacy guarantee
whose unit of protection is also the person; user-level differential privacy is
the granularity that matches the threat model, because it is subject
participation, not window participation, that a twin could disclose.

The results place that requirement in tension with feasibility. User-level
privacy is precisely the granularity that this study finds infeasible at strong
budgets for a clinical-scale cohort: at $\varepsilon = 0.5$ the four-class model
is statistically indistinguishable from chance. Sample-level privacy, which
remains useful at the same budget, protects the wrong unit and additionally
eliminates the minority class. The twin setting therefore cannot simply adopt the
mechanism with the better aggregate curve; it inherits the harder of the two
problems. This tension is, in our view, the strongest motivation for the
directions the programme takes next: larger cohorts, at which user-level
amplification becomes available and the strong-budget regime becomes reachable,
and personalised federated schemes that may retain subject-specific signal ---
the signal federation attenuates --- without transmitting it. The present paper
establishes the measurement framework and the mechanism-level findings on which
that twin-specific work will build.
